\documentclass[11pt,a4paper]{article}
% Shared preamble for RuPhO English problem statements (match T1 2026 style).
\usepackage[margin=1in]{geometry}
\usepackage{amsmath,amssymb}
\usepackage{graphicx}
\usepackage{float}
\usepackage{enumitem}
\usepackage{microtype}
\usepackage{caption}
\usepackage{tabularx}
\usepackage{hyperref}

\captionsetup{font=small,labelfont=bf,skip=6pt}

\setlist[itemize]{leftmargin=1.2cm}
\setlist[enumerate]{leftmargin=1.2cm}

% One outer box per subproblem: [id | statement | points] (no inner rules)
\newcommand{\problembox}[3]{%
  \par\addvspace{0.42em}%
  \noindent
  \setlength{\fboxsep}{7.5pt}%
  \setlength{\fboxrule}{0.95pt}%
  \fbox{%
    \begin{minipage}{\dimexpr\linewidth-2\fboxsep-2\fboxrule\relax}
    \begin{tabularx}{\linewidth}{@{}>{\raggedright\arraybackslash\bfseries}p{2.1em} @{\hspace{0.9em}} >{\raggedright\arraybackslash}X @{\hspace{0.9em}} >{\raggedleft\arraybackslash\bfseries}p{2.4em}@{}}
    #1 & #2 & #3
    \end{tabularx}
    \end{minipage}%
  }%
  \par\addvspace{0.32em}%
}

\title{\textbf{Rocket engine}}
\author{\normalsize Y17 (2017) --- English translation}
\date{}
\begin{document}
\maketitle

In this problem we will study the flow of gas in the nozzle of a rocket engine,
and then the flight of the rocket in airspace and in vacuum.

To begin with, let us recall several fundamental relationships that describe the
behavior of flowing liquids or gases.

\section*{Part A. Dynamic relations for gas (1.0 points)}

Let us consider a small cylinder of matter with density $\rho$ (see figure),
moving under the influence of a pressure difference.
The movement is stationary and one-dimensional.

\begin{figure}[H]
    \centering
    \includegraphics[width=0.72\textwidth]{T5_EN_assets/fig1.jpg}
    \caption{Figure 1.}
\end{figure}

\problembox{A1}{Write down the equation of motion of the cylinder.
Determine the pressure gradient $dp/dx$.
Express your answer in terms of $\rho$, $v$ and $dv/dx$.}{0.50}

\problembox{A2}{Determine the speed of sound $c$ in the gas.
Express your answer in terms of $p$, $\rho$ and the adiabatic exponent $\gamma$.}{0.50}

\section*{Part B. Mach numbers and nozzle profile (2.0 points)}

To describe gas flows, it is convenient to introduce the Mach number $M=v/c$,
which characterizes the ratio of the local gas flow velocity to the speed of sound.

We now move on to a qualitative description of the rocket engine nozzle.
In a rocket engine, fuel burns in the combustion chamber and enters the nozzle.
The speed of combustion products increases throughout the nozzle.

Solve the following points of the problem under the following assumptions:
the gas is considered ideal;
the gas flow is isentropic (constant entropy, no friction or dissipative losses)
and adiabatic (no heat exchange);
the gas flow is stationary and one-dimensional, that is, at any fixed point of the
nozzle all flow parameters are constant in time and change only along the nozzle axis,
at all points of a given cross section the flow parameters are the same,
and the gas velocity vector is everywhere parallel to the symmetry axis of the nozzle;
the gas mass flow rate is the same in all cross sections of the flow;
the influence of all external forces and fields (including gravitational) is negligible;
the symmetry axis of the nozzle is the spatial coordinate $x$.

\problembox{B1}{Let us denote the cross-sectional area of the nozzle as $A$.
Determine the value of $dA/dx$.
Express your answer in terms of $A$, $v$, $c$ and $dv/dx$.}{1.00}

\problembox{B2}{Qualitatively construct the profile of the nozzle section $A(x)$
when the gas velocity at the outlet is $v{<}c$.
We can assume that at the entrance to the nozzle the gas velocity is negligible.
Indicate special points (if any).}{0.50}

\problembox{B3}{Qualitatively construct the profile of the nozzle section $A(x)$
when the gas velocity at the outlet is $v{>}c$.
We can assume that at the entrance to the nozzle the gas velocity is negligible.
Indicate special points (if any).}{0.50}

\section*{Part C: Analysis of Nozzle Flow (5.0 points)}

Let us consider in more detail the process of gas flow in the nozzle.
Let the gas pressure $p_c$, temperature $T_c$ and density $\rho_c$ at the nozzle inlet.
Gas adiabatic index $\gamma$.

\problembox{C1}{When making a rocket engine, the desired pressure at the nozzle
exit $p$ is set. Find the gas flow rate $v$ from the nozzle as a function of the
ratio $p/p_c$. Express your answer in terms of $p_c$, $\rho_c$, $\gamma$ and
the ratio $p/p_c$.}{2.00}

\problembox{C2}{At what value of $p'$ does the outflow velocity from the nozzle
reach its maximum value? What is this maximum speed $v_\text{max}$?
Express your answers in terms of $p_c$, $\rho_c$ and $\gamma$.}{0.40}

\problembox{C3}{How does the nozzle outlet area $A$ depend on the desired pressure
at the nozzle outlet $p$ (with fixed parameters at the nozzle inlet)?
Find the function $A(p/p_c)$, assuming that the mass flow rate in each section
of the nozzle is constant and equal to $\lambda$.
Express your answer in terms of $\lambda$, $p_c$, $\rho_c$, $\gamma$ and
the ratio $p/p_c$.}{1.10}

\problembox{C4}{Find the pressure ratio $p_t/p_c$ at the singular points for
points B2 and B3. Express your answer in terms of the gas adiabatic index $\gamma$.
Such pressures are called critical.}{1.50}

\problembox{C5}{For the nozzle shape defined in step B3, determine the dependence
of the value $v/c$ on the value $p/p_c$. Express your answer in terms of $\gamma$
and $p/p_c$. Construct qualitative graphs of the dependence of the gas flow
velocity $v/c$ along the nozzle for the cases when the outlet pressure is greater
than the critical one ($p/p_c>p_t/p_c$) and when the outlet pressure is less than
the critical one ($p/p_c{<}p_t/p_c$).}{0.70}

\problembox{C6}{For the nozzle shape defined in point B3, the gas mass flow rate
at the nozzle exit can be represented as
$\lambda = C(A_\text{out},\gamma,p_c,\rho_c)\cdot f(p/p_c)$,
where $C$ is a constant that depends only on $A_\text{out}$, $\gamma$, $p_c$,
and $\rho_c$, and $f(p/p_c)$ is a function that depends only on the ratio $p/p_c$
and $\gamma$. Determine the function $f(p/p_c)$. Express your answer in terms of
$p/p_c$ and $\gamma$. Construct a qualitative graph of the mass flow rate $\lambda$
as a function of the pressure at the nozzle exit ($p/p_c\in[0;1]$). If the graph
has special points, indicate them on the graph and obtain the corresponding
analytical expressions for $p/p_c$.}{0.80}

\section*{Part D. Momentum acquired by a rocket (0.5 points)}

Let the rocket now fly in the atmosphere. Atmospheric pressure $p_0$, gas pressure
at the nozzle outlet $p$, nozzle area $A$, gas mass flow $\lambda$,
gas flow rate from the nozzle $v$.

\problembox{D1}{Find the specific impulse acquired by the rocket while flying
in the atmosphere. Express your answer in terms of $p_0$, $p$, $A$ and $v$.
Specific impulse is the additional impulse acquired by a rocket when releasing
a unit mass of fuel.}{0.50}

\vspace{1em}
\noindent\textit{Source:} \url{https://pho.rs/p/3040}

\end{document}
